\chapter[Conclusion]{Conclusion and Future Work}
\label{chap:conclusion}

\section{Answer to the Research Question}

The question this project set out to answer was not whether retrieval-augmented
generation can answer questions about Regulation (EU) No 376/2014. It was
\emph{where such a system fails, and which condition observable at runtime
justifies escalating instead of answering}. The distinction matters because a
single accuracy figure conceals the first question entirely and cannot address
the second at all.

The answer has three parts, none of which was the answer anticipated at the
outset.

\subsection{Retrieval is not where this system fails}

On the frozen benchmark, hybrid retrieval with cross-encoder reranking places
acceptable evidence in the top five results for 89.3\% of held-out questions and
in the top ten for 94.6\%. The candidate pool at depth 20 contains acceptable
evidence for 96.4\%. Retrieval was the natural suspect, and the measurements do not support it:
by the time the generator sees its context, the evidence needed to answer is
usually present.

What the retrieval work established is narrower than a percentage and bears
more directly on the policy. The lexical and dense channels fail on different questions, so fusion addresses a coverage gap rather than tuning an existing
configuration; reranking changes the order of a fixed
pool without enlarging it, so anything absent at depth 20 is unrecoverable
downstream and must be handled as a distinct condition rather than as a bad
answer. Those two facts shape the escalation policy more than the recall figures
do.

\subsection{The failure is citation precision, and it is invisible to
conventional measures}

The historical baseline, re-scored on frozen labels and regenerated over the
whole benchmark, locates the failure more precisely. Retrieval-augmented
generation identifies the correct legal parent for 62.5\% of questions but
reproduces the complete citation, sub-paragraph included, for only 44.3\%. Those
eighteen points are sixteen answers that retrieved the right article, read
fluently, and cited the wrong paragraph of it.

The original evaluation could not see this. Its citation measure awarded credit
for substring containment and half credit for naming the article number alone, so
a prediction citing Article 16(1) where the provision is Article 16(6) was
recorded as correct. Under the corrected evaluator, few-shot prompting falls from
0.286 to 0.159, because naming a real article with the wrong sub-paragraph was
precisely its characteristic error, and precisely what the old measure rewarded.

The ranking of the three historical systems is unchanged. The magnitudes are
not, and neither is the population they were measured on: an earlier draft
reported this comparison over 73 records, because 25 of the 98 had been reworded
after the stored predictions were made. Regenerating those predictions against the
frozen wording shows the exclusion was not neutral. On the 25 records that had
been dropped, retrieval-augmented generation reaches 0.360 at parent level against
0.730 on the 73 that were kept, with non-overlapping intervals, while few-shot
moves the other way. The questions an audit rewrites for vagueness are harder for
a retriever and easier for a model answering from memory. All three statements are
reported together.

\subsection{The escalation condition, and the honest limit on it}

Two conditions observable at runtime were examined. One works and needs no
calibration at all; the other cannot yet be calibrated on data of this size.

\textbf{Quote verification is a usable gate.} Whether a quoted span appears
contiguously in the source it cites is decidable without gold labels, without
model judgement and without a threshold. It has no false positives by
construction. On the held-out partition it rejected 7 of 56 answers, and every
rejection was a genuine defect: elided quotes that deleted conditions from legal
tests, a paraphrase presented as quotation, and a fabricated source identifier one
character removed from a real one.

\textbf{Retrieval-uncertainty escalation cannot be calibrated on data of this
size, and no threshold for it is reported.} With nine error events on the development
partition, every candidate rule's confidence interval on precision contains or
nearly contains the base error rate. The most promising signal --- disagreement
between rank fusion and the reranker about which result belongs first --- reaches
0.42 precision at 0.89 recall, with an interval of [0.23, 0.64] against a base
rate of 0.21. That is a direction to investigate rather than a threshold. Demonstrating the
observed effect would require roughly 178 questions against a benchmark of 98.

It is reported as a negative result with a stated sample-size requirement
rather than as a number fitted to nine events.

\section{What the Project Established That It Did Not Set Out to Find}

Four results emerged from the work rather than from its design, and they are the
ones most likely to transfer beyond this corpus.

\subsection{Verification and instruction are not substitutes}

The same negative result appears three times.

Instructed to quote verbatim, the generator elided material. Instructed
explicitly not to elide, it stopped marking omissions with an ellipsis and began
splicing fragments silently instead --- a failure a syntactic check for ellipsis
characters would have passed, and which only contiguity matching rejected. The
prompt rule that eliminated elision on the development partition did not
eliminate it on unseen questions, where it recurred in 4 of 56 answers. And a
response the provider's strict structured-output mode had already accepted
carried a source identifier the application had never supplied.

In each case an instruction shifted the failure rather than removing it, and in
each case only a check against the evidence detected what remained.
\textbf{A rule stated to a generator is a request; a rule checked against the
evidence is a guarantee.} It is the project's central practical claim, and it was reached by trying the
cheaper option first and measuring the outcome.

\subsection{Abstention cannot be delegated to the generator}

Across 140 generated answers on both partitions, the model abstained zero times.
The output schema offered four abstention reasons, the prompt instructed their
use, and strict structured output guaranteed that returning one would be
well-formed and cost nothing. It abstained on none of the eighteen answers that
cited the wrong provision, none of the answers whose quotes were not grounded,
and not on the answer that invented an identifier.

Selective-prediction research generally assumes a model can be induced to signal
when it should not answer. On this task, with this model, the explicit and freely
available channel for doing so went entirely unused. Abstention here is therefore
not a behaviour to elicit but a decision to impose from outside, on evidence the
system can check.

\subsection{The two failure modes are independent}

On the development partition, nine answers selected the wrong provision and two
failed quote verification. \emph{No answer did both.} Quote verification has
precision 0.00 and recall 0.00 for predicting evidence-selection errors.

The consequence is architectural. Misreporting what a provision says is
detectable at runtime; selecting the wrong provision is not, and must be
predicted from retrieval signals. A verifier addressing only the first would have
passed all nine errors of the second. Two layers are required, and the
requirement is measured rather than assumed.

\subsection{A benchmark's labels can be internally valid and externally
unreadable}

Twenty-two of the 98 frozen records carry a citation label that cannot be
resolved to a single provision from the string alone: twelve omit which of four
regulations they mean, and ten cannot be parsed into a provision at all. Every
one of them resolves correctly inside the pipeline, which carries the parent
identifier alongside the display string. The defect is invisible to every
internal metric by construction and immediately visible to a practitioner holding
only what the user is shown.

It was found because reviewers were asked, not because anything was measured.
Six of them wrote variants of \emph{\enquote{GM 3.1 = ?}} in a free-text box. The
measurement came afterwards and showed the problem to be roughly four times
larger than the comments suggested.

\section{Contributions}

\begin{enumerate}
  \item \textbf{A deterministic, hierarchy-aware corpus} of 109 legal parents and
    1{,}535 retrieval children with byte-identical rebuild and zero validation
    issues, together with explicit lineage from the historical flat chunks.

  \item \textbf{An independently reviewed benchmark of 98 evidence-linked
    questions}, split by legal parent with zero leakage, frozen under a
    pre-registered rule with the review's disagreement reported rather than
    resolved away. Fleiss' $\kappa$ of 0.250 across four practitioners is itself
    a finding: a single gold label was never as certain as a single accuracy
    figure implied.

  \item \textbf{A citation evaluator that resolves references against the corpus}
    rather than matching strings, reported at two levels of strictness, with the
    scorer it replaces retained and its errors pinned by tests so the change is
    auditable rather than asserted.

  \item \textbf{A deterministic verification layer} that checks source
    admissibility and quote contiguity without gold labels or model judgement,
    and the measurement showing why it must sit alongside, not instead of, a
    check on evidence selection.

  \item \textbf{An escalation policy specified against runtime-observable
    conditions}, with the one component that could be calibrated left explicitly
    uncalibrated and the required sample size stated.

  \item \textbf{A methodological record} of 45 append-only decisions and a
    disclosure log of every AI-assisted artefact with what was verified and what
    remained outstanding, sufficient for the work to be reconstructed and
    challenged.
\end{enumerate}

\section{Limitations}

\textbf{The benchmark is small and synthetic in origin.} Ninety-eight records,
33 of which were independently reviewed, is enough to characterise failure modes
and not enough to calibrate a policy against them. Its questions were generated
by a language model and may reflect that model's style.

\textbf{Reviewer agreement is fair, not strong.} A $\kappa$ of 0.250 bounds how
precisely any figure computed against these labels can be read. This is reported
as a property of the task rather than a defect of the panel, but it is a real
limit on interpretation.

\textbf{Semantic support is unmeasured.} Every verification result in this work
is structural. That a claim points at supplied evidence and reproduces its wording
contiguously does not establish that the evidence supports the claim. The 83.9\%
held-out evidence-selection figure counts whether the cited source is among the
acceptable ones, not whether the sentence built on it is sound.

\textbf{One model, one language, two configurations, one of them development-only.}
Generation used a single model at a single temperature. The neighbour-enriched
context variant was evaluated on the development partition and not adopted: it
costs 79\% more input tokens and moves evidence selection and byte-exact quoting
by one record each, which 42 questions cannot separate from noise. Retrieval is English-first, and the reranker was
chosen for local execution and English passage ranking. Nothing here speaks to
multilingual behaviour, which matters for a regulation applied across the Union.

\textbf{The LoRA arm of the comparison covers 73 of 98 records.} Few-shot and
RAG were regenerated over the frozen wording and cover all 98; LoRA was not
re-run, so its stored predictions are scored on the 73 records whose wording did
not change. That subset is measurably easier for retrieval-augmented systems, so
any comparison against LoRA flatters RAG rather than the reverse.

\textbf{Equal-weight fusion is deliberately suboptimal.} Convex combination of
normalised scores outperforms reciprocal-rank fusion when its weights are tuned.
Tuning them would have required fitting on data that also grounds the reported
results, so the retrieval figures are a conservative estimate of what the
architecture can reach.

\textbf{This is a research prototype.} It does not replace qualified
interpretation of aviation law, and nothing in this work should be read as
suggesting it could.

\section{Future Work}

\textbf{Extend the benchmark to the size the calibration requires.} The power
analysis gives the target directly: roughly 178 questions to demonstrate the
escalation precision actually observed, against 98 today. This is the single
change that would convert the escalation policy from specified to calibrated, and
it is a data-collection task rather than a research one.

\textbf{Build the cross-reference subset before building an agent.} Every record
in the frozen benchmark has its gold evidence inside a single legal parent; none
requires combining provisions from different parts of the instrument. A benchmark
with no multi-hop questions cannot exhibit a multi-hop system's benefit, so
agentic retrieval remains unevaluable here regardless of its merits elsewhere.
The prerequisite is a subset of questions whose evidence genuinely spans parents,
constructed and reviewed to the same standard as the present set.

\textbf{Measure semantic support, not only structural grounding.} The natural
next layer is an entailment check between each claim and the evidence it cites,
which would close the gap between \enquote{quotes the source correctly} and
\enquote{is supported by the source}. Given what this project found about
delegating judgement to a generator, such a check should be validated against
human ratings before being trusted, and its own agreement with those ratings
reported.

\textbf{Repair the citation labels, then re-measure.} Twenty-two records carry
labels a reader cannot resolve. Rewriting them to name their regulation and
provision unambiguously is mechanical, and it would allow the citation-legibility
finding to be stated as a repaired defect rather than only as a diagnosis.

\textbf{Test whether the escalation signals transfer across model capability.}
The generator was deliberately held fixed at the model used for the historical
baselines, so that the comparison isolated scaffolding from capability. Whether a
policy calibrated on a weaker generator still holds for a stronger one is an open
and practically important question, and running the existing pipeline against a
current model would answer it cheaply.

\section{Closing Remark}

The system this work produced answers questions about occurrence reporting with
retrieved evidence, checks its own quotations against that evidence, and refuses
to answer when the check fails. It is not a system that is right more often than
its predecessor by a large margin. It is a system whose errors can be detected, and detectability is the property
this domain requires.

Regulatory questions are asked by people who will act on the answer. For them,
a confident wrong citation is worse than a referral to a human, and an answer
that cannot be audited is worse than no answer. The finding that bears most directly on deployment is that the generator will
not signal when it is in that position ---
140 opportunities, zero abstentions --- and that the system must therefore be built
to determine it independently.
